% Generated draft: requires mathematical and source review.
\documentclass[11pt]{article}
\usepackage[margin=25mm]{geometry}
\usepackage{fontspec}
\setmainfont{Latin Modern Roman}
\usepackage{amsmath,amssymb,mathtools,unicode-math}
\setmathfont{Latin Modern Math}
\usepackage{enumitem,needspace,xurl,hyperref}
\setlength{\parindent}{0pt}
\setlength{\parskip}{5pt}
\setlength{\emergencystretch}{3em}
\newcommand{\N}{\mathbb{N}}
\newcommand{\Z}{\mathbb{Z}}
\newcommand{\Q}{\mathbb{Q}}
\newcommand{\R}{\mathbb{R}}
\newcommand{\C}{\mathbb{C}}
\newcommand{\F}{\mathbb{F}}
\newcommand{\A}{\mathbb{A}}
\newcommand{\PP}{\mathbb{P}}
\newcommand{\E}{\mathbb{E}}
\newcommand{\eps}{\varepsilon}
\newcommand{\dd}{\,\mathrm{d}}
\begin{document}
\section*{UnsolvedMath ID: 701068}
\textbf{Source problem code:} catalogue entry
\subsection{Definitions and mathematical statement}
A group is called a \emph{Chernikov group} if it has a normal subgroup of finite index isomorphic to a direct product of finitely many quasicyclic groups \(C_{p^\infty}\), where the primes may vary and the empty product is allowed. Here \(C_{p^\infty}\) is the group of all complex roots of unity whose orders are powers of the prime \(p\). In accordance with the source, for a group property \(\sigma\), a group is \emph{quasi-\(\sigma\)} if it does not have property \(\sigma\), but every proper subgroup has property \(\sigma\). Thus a quasi-Chernikov group is not Chernikov and every proper subgroup is Chernikov; in particular, it is infinite.

A group is \emph{simple} if it is nontrivial and its only normal subgroups are the identity subgroup and the whole group. An involution is an element of order exactly \(2\). Write \(b^g=g^{-1}bg\), and let \(\langle a,b^g\rangle\) denote the subgroup generated by \(a\) and \(b^g\). An infinite simple group \(G\) is a \emph{monster of the third kind} if, for every \(a,b\in G\setminus\{1\}\) with \(a^2\ne1\) or \(b^2\ne1\), the set
\[
\{g\in G:\langle a,b^g\rangle=G\}
\]
is infinite. The infinitude here concerns the conjugating elements \(g\), not necessarily distinct conjugates of \(b\).

\textbf{Question (A. I. Sozutov, Kourovka Notebook 14.83).} Is every simple quasi-Chernikov group a monster of the third kind? Equivalently, does every simple group \(G\) that is not Chernikov and whose proper subgroups are all Chernikov satisfy
\[
\forall a,b\in G\setminus\{1\},\qquad
(a^2\ne1\ \lor\ b^2\ne1)
\ \Longrightarrow\
\bigl|\{g\in G:\langle a,g^{-1}bg\rangle=G\}\bigr|\text{ is infinite}?
\]
The supplied source reports an affirmative answer for simple quasi-finite groups, meaning infinite simple groups all of whose proper subgroups are finite, and attributes this special case to Sozutov's 1997 paper. This is a supplied literature assertion, not an independently verified result here. The source also requests comparison with V. P. Shunkov's definitions in Archive problems 6.63 and 6.64; those definitions are not supplied.
\subsection{Short English statement}
Must every infinite simple group that is not Chernikov, but has only Chernikov proper subgroups, have the following property: for any two nonidentity elements, at least one of which is not an involution, infinitely many conjugating elements make the first element and a conjugate of the second generate the whole group? The supplied record labels this problem open and reports the quasi-finite case as known.
\subsection{Sources}
\begin{itemize}
\item Local UnsolvedMath registry: permanent ID \(701068\), exact source problem code \texttt{catalogue entry}, title \emph{Kourovka Notebook Problem 14.83}. The supplied record labels the problem \emph{open}; this label is not independently verified here.
\item A. I. Sozutov, problem 14.83, \emph{The Kourovka Notebook (Unsolved Problems in Group Theory)}, 14th issue (1999). The supplied record identifies the compilation as the 21st issue edition (2026), arXiv:1401.0300v48, and names E. I. Khukhro and V. D. Mazurov as editors of the Issues 1--21 collection.
\item A. I. Sozutov, \emph{Algebra and Logic} \textbf{36}, no. 5 (1997), 336--348. The source cites this paper for the affirmative quasi-finite case; no article title or link is supplied.
\item V. P. Shunkov, definitions in the Kourovka Notebook Archive, problems 6.63 and 6.64, cited for comparison in the supplied statement. Their contents are not included in the record.
\end{itemize}
\end{document}
